\documentclass[10pt]{article}
\usepackage[letterpaper,margin=0.7in]{geometry}
\usepackage{enumitem}
\usepackage{titlesec}
\usepackage[hidelinks]{hyperref}
\usepackage{xcolor}
\usepackage{ragged2e}
\usepackage{fontawesome5}
\usepackage[default]{raleway}
\usepackage{multicol}

% Colors
\definecolor{accent}{RGB}{10,60,90}
\definecolor{secondary}{RGB}{30,30,30}

% Font settings
\renewcommand{\familydefault}{\sfdefault}
\newcommand{\mytitle}[1]{\color{accent}\normalfont\large\bfseries{#1}}

% Section formatting
\titleformat{\section}{\mytitle}{}{0em}{}[{\titlerule[0.8pt]}]
\titlespacing{\section}{0pt}{9pt}{5pt}

% Itemize styling
\setlist[itemize]{leftmargin=*,nosep,after=\vspace{3pt}}
\setlist[itemize]{label=\color{accent}\scriptsize\raisebox{0.4ex}{\textbullet}}

% Custom commands
\newcommand{\tech}[1]{\textcolor{secondary}{\textbf{#1}}}

\newcommand{\experience}[4]{%
  \vspace{4pt}
  \begin{tabular*}{\textwidth}{@{}l@{\extracolsep{\fill}}r@{}}
    \textbf{#1} & \textit{#2} \\
    \textit{\small#3} & \textit{\small#4}
  \end{tabular*}
  \vspace{-4pt}
}

\begin{document}

% ===== Header =====
\begin{center}
    {\Huge \textbf{JOANNE HOAR, M.Sc.}} \\[4pt]
    {\Large \tech{Software Developer \textbullet\ Technology Educator \textbullet\ AI Practitioner}} \\[6pt]
    \faIcon{at}~\href{mailto:joanne.hoar@gmail.com}{joanne.hoar@gmail.com}
    \hfill
    \faIcon{phone}~+1\,403\,854\,8433
    \hfill
    \faIcon{linkedin}~\href{https://linkedin.com/in/joanne-hoar/}{linkedin.com/in/joanne-hoar}
    \hfill
    \faIcon{github}~\href{https://github.com/xcodegirl}{xcodegirl}
    \\[2pt]
    \textit{Hanna, Alberta, Canada}~\textbullet~\textit{Remote}
\end{center}

% ===== Summary =====
\section*{Summary}

Canadian Software Developer and Technology Educator with a Master of Science in Computer Science (University of Calgary) and more than twenty years of experience building technical systems and teaching. Currently an active online instructor at Robertson College, where I teach and maintain curriculum across twelve programming and data subjects on Brightspace/D2L. I use AI tools daily across content development, code generation, CLI workflows, and process automation, and I teach students how to use AI responsibly, as a tool rather than a solution. I have developed and actively deliver curriculum on AI ethics, responsible use, and academic integrity as it applies to AI-generated content.

% ===== AI Expertise =====
\section*{AI Expertise}

\textbf{Daily practice:}
\begin{itemize}
    \item Generative AI for content development: drafting course materials, case studies, scaffolded exercises, and code examples, always with accuracy review and editorial refinement
    \item Prompt engineering for structured output: system prompts, iterative refinement, chain-of-thought elicitation, output format and style control
    \item AI-assisted coding: \tech{GitHub Copilot}, \tech{Claude}, \tech{ChatGPT}: for scaffolding, debugging, refactoring, and code review
    \item CLI and workflow automation: using AI to interface with \tech{Azure}, \tech{GitHub}, and shell tooling
    \item Multi-model practice: Claude, ChatGPT, GitHub Copilot, DeepSeek, Recraft AI, with different tools for different tasks
\end{itemize}

\textbf{What I teach about AI:}
\begin{itemize}
    \item Responsible AI use as professional practice: practitioners are ultimately accountable for everything they produce, with or without AI assistance
    \item \textbf{Three-level AI assessment framework} (adapted from Perkins et al., 2024): from no-AI to AI-assisted task completion with mandatory human evaluation, applied directly to programming assignments
    \item Citation of AI-generated code: inline comment standards, README disclosure, terms of service compliance, with real code examples across seven scenarios
    \item Copyright, bias, hallucination detection, and institutional policy: applied to software contexts, not just essay writing
    \item AI tools as a first draft requiring editorial judgment, not a finished product
\end{itemize}

\textbf{AI and assessment design} \textit{\small(from Robertson's academic integrity training and the Academic Integrity Inter-Institutional Meeting, AIIIM 2026)}\textbf{:}
\begin{itemize}
    \item ``Authentic'' learning and assessment: assignments grounded in real, local contexts, where generic AI output cannot succeed without genuine human judgment
    \item Scaffolded assignments with checkpoints: milestone grading and visible work in progress rather than a single final deliverable
    \item Developing professional judgment of AI output: students audit, correct, and debug flawed, biased, or hallucinated AI drafts
    \item Peer review and peer grading of AI-assisted work against concrete rubrics
\end{itemize}

% ===== Teaching Experience =====
\section*{Teaching \& Course Development}

\experience{Technology Instructor}{Jan 2025 -- Present}{Robertson College \normalfont\small\textit{-- Online, Brightspace / D2L (Full Stack Development \& Data Analytics programs)}}{}
\begin{itemize}
    \item Maintain and deliver all course materials for \textbf{12 subjects}, with \textbf{20 course offerings} taught since January 2025; contribute to course revisions and create supplemental materials: sample code repositories, reference sheets, recorded walkthroughs, and open-source GitHub archives
    \item Teach students to use AI responsibly in software development; deliver academic integrity curriculum on AI citation, policy, and ethics; developed a dedicated slide deck actively used each term
    \item Completed Robertson's academic integrity training; participated in the Academic Integrity Inter-Institutional Meeting (AIIIM 2026)
    \item Served as \textbf{contracted Subject Matter Expert} for the Full Stack Web Developer program: wrote its sections of the cross-program Tech Programs FAQ for Robertson's internal Program Knowledge course, then built and presented slides on the program to about forty Robertson team members at an internal tech training session with live Q\&A (March 2025)
    \item Invited back for future sessions and ongoing FSD SME consultation following the engagement
    \item Conduct 3+ synchronous live coding sessions per week; use real-time debugging as deliberate instructional modelling of professional problem-solving
    \item Apply Universal Design for Learning throughout: captions on recordings, structured documents, hardware-agnostic assignments
    \item Student satisfaction: \textbf{4.5\,/\,5.0} teaching career-changing adult learners; feedback highlights real-world context, approachability, and responsiveness
    \item[] \tech{Brightspace/D2L \textbullet\ Angular \textbullet\ React \textbullet\ ASP.NET Core \textbullet\ C\# \textbullet\ Java \textbullet\ Git \textbullet\ Azure \textbullet\ SQL \textbullet\ Python}
\end{itemize}

\experience{Instructor, Intermediate C++ Programming}{2005 -- 2008}{University of Calgary, Department of Continuing Education \normalfont\small\textit{-- Adult learners}}{}
\begin{itemize}
    \item Designed and delivered complete curriculum for adult learners: outcomes, lectures, sample code, and assessments from the ground up
\end{itemize}

\experience{Research / Teaching Assistant}{2001 -- 2004}{VLab / AI Research Hub, University of Calgary}{}
\begin{itemize}
    \item TA and lab instructor: \textit{Introduction to Programming}, \textit{Computing for Non-Majors}, and \textit{Operating Systems}
    \item Computing for Non-Majors required translating technical concepts for a completely non-technical audience
\end{itemize}

% ===== Industry Experience =====
\section*{Industry Experience}

\experience{Senior Software Developer}{2022 -- 2025}{bitHeads Inc. \normalfont\small\textit{-- brainCloud BaaS SDK}}{}
\begin{itemize}
    \item Authored extensive developer-facing documentation, onboarding guides, and sample projects; hired in part for educational background
    \item Provided one-on-one client instruction across seven UI frameworks and twelve platforms
\end{itemize}

\experience{Co-founder / Lead Developer}{2005 -- 2022}{Silicon Hanna Inc. \normalfont\small\textit{(self-employed; founded as SuRJE Software Solutions Inc.)}}{}
\begin{itemize}
    \item 17-year software development and consulting practice spanning mobile, web, AR, and games
    \item Early AI/ML research roots: M.Sc. in agent-based systems, swarm intelligence, and evolutionary computation, the academic precursors to today's machine learning landscape
    \item Published 11 mobile titles (iOS and Android); explored novel SDKs and produced learning materials throughout
\end{itemize}

\experience{Tools Programmer}{2008 -- 2011}{Fekete Associates Inc. (S\&P Global)}{}
\begin{itemize}
    \item Designed and delivered a voluntary \textbf{internal graphics training course} for colleagues from beginners to senior developers; built multi-level materials with a hands-on capstone
\end{itemize}

\experience{Intermediate Software Engineer}{2005 -- 2008}{BJ Pipeline Inspection Services Inc. (Baker Hughes)}{}
\begin{itemize}
    \item Implemented \textbf{neural network algorithms} for automated anomaly detection in pipeline inspection data: applied machine learning in industry well before the current AI wave
\end{itemize}

% ===== Education =====
\section*{Education}
\begin{itemize}
    \item \tech{M.Sc. Computer Science} -- University of Calgary \hfill 2004 \\
    \textit{Thesis: Agent-based Development of Natural Transportation Networks} \\
    \textit{Graduate coursework: Human-Computer Interaction, Artificial Intelligence, Computer Graphics}
    \item \tech{B.Sc. Computer Science} -- University of Calgary \hfill 2001
\end{itemize}

% ===== Honours =====
\section*{Honours \& Awards}
\begin{itemize}[nosep,leftmargin=*]
    \item Departmental Nomination: Governor General's Gold Medal, Faculty of Graduate Studies \hfill 2005
    \item Postgraduate Scholarship NSERC D3, Natural Sciences and Engineering Research Council of Canada \hfill 2005
    \item Graduate Student Scholarship, Alberta Learning \hfill 2004
    \item Best Student Paper Award (Co-author), Australian Conference on Artificial Life (ACAL) \hfill 2003
    \item CPSC Departmental Research Award, University of Calgary \hfill 2003
\end{itemize}

% ===== Publications =====
\section*{Selected Publications \normalfont\small\textit{(published under maiden name Penner)}}
\begin{itemize}[nosep,leftmargin=*]
    \item \textbf{Penner, J.}, Hoar, R., Jacob, C. ``Bacterial Chemotaxis in Silico,'' \textit{Australian Conference on Artificial Life (ACAL)}, 2003 \hfill \textit{Best Student Paper}
    \item Hoar, R., \textbf{Penner, J.}, Jacob, C. ``Transcription and Evolution of a Virtual Bacteria Culture,'' \textit{IEEE Congress on Evolutionary Computation (CEC)}, 2003
    \item Hoar, R., \textbf{Penner, J.} ``The Application of Artificial Intelligence to Transportation System Design,'' \textit{XRDS: Crossroads, ACM}, 2003
    \item \textbf{Penner, J.}, Hoar, R., Jacob, C. ``Swarm-Based Traffic Simulation with Evolutionary Traffic Light Adaptation,'' \textit{IASTED ASM}, 2002
    \item Hoar, R., \textbf{Penner, J.}, Jacob, C. ``Evolutionary Swarm Traffic: If Ant Roads Had Traffic Lights,'' \textit{IEEE CEC}, 2002
\end{itemize}

% ===== Skills =====
\section*{Technical Skills}
\textbf{AI Tools:} ChatGPT, Claude, GitHub Copilot, DeepSeek, Recraft AI, used responsibly in curriculum, content, and coding workflows \\
\textbf{LMS:} Brightspace / D2L (daily), GitHub (open-source course repos) \\
\textbf{Languages:} C++, C\#, JavaScript, TypeScript, Java, Python, PHP, SQL \\
\textbf{Frameworks:} Angular, React, ASP.NET Core, Node.js, .NET \\
\textbf{Content \& Authoring:} Markdown, \LaTeX, HTML/CSS, open-source tooling \\
\textbf{Human Languages:} English (native), French (intermediate), Spanish (rudimentary)

% ===== Work Samples =====
\section*{Work Samples}
\faIcon{github}~\href{https://github.com/xcodegirl/career/blob/main/xcodegirl-work-samples.md}{github.com/xcodegirl/career (xcodegirl-work-samples.md)} \\
Sample lesson plans and slide decks available upon request.

\end{document}
